\documentclass[11pt,a4paper]{article}
\usepackage[margin=2.4cm]{geometry}
\usepackage{amsmath,amssymb,booktabs,hyperref}
\newcommand{\code}[1]{\texttt{#1}}
\newcommand{\dt}{\Delta t}
\title{The Exner pressure as an additional prognostic variable\\in the frozen split (dry and moist)}
\author{MultiRate notes}
\date{\today}
\begin{document}
\sloppy
\maketitle

\begin{abstract}
Following the pressure-update idea of J\"ahn and Knoth (poster ``Compressible modeling of a cloudy
atmosphere using a general pressure evolution equation''), the Exner pressure $\pi$ is carried as an
additional prognostic variable next to the conservative thermodynamic variable. The pressure gradient
is written as the product $c_p\hat\Theta\,\nabla\pi$, the frozen split of the MIS method integrates
$\pi$ with its own evolution equation in the fast substeps, and $\pi$ is recomputed from the equation of
state after every MIS step (not at the internal stages). The dry formulation (\code{FV.ExnerModel},
\code{src/FV/exner.jl}) and the moist formulation with reversible thermodynamics
(\code{FV.MoistExnerModel}, \code{src/FV/moist.jl}) are described for the flat C-grid finite volumes,
together with the variant that carries the pressure $p$ itself (option \code{pvar = :p}, the standard
pressure equation in the form that follows from the conservative variables) and tests: dry warm
bubble, gravity-wave channel, moist bubble of Bryan and Fritsch (2002). A moist Exner pressure (moist
exponent) is discussed and not recommended.
\end{abstract}

\section{Idea}
The thermodynamic variable $\Theta$ (dry: $\rho\theta$; moist: $\rho\theta_\rho$) is transported
conservatively. The pressure is a function of the state, $p=P(\rho,\Theta,\dots)$; instead of evaluating
$P$ in every fast substep, the Exner pressure obeys
\[
\partial_t\pi=\frac{\partial\pi}{\partial\rho}\,\partial_t\rho+\frac{\partial\pi}{\partial\Theta}\,\partial_t\Theta+\dots,
\]
which keeps $\pi=\pi(\rho,\Theta,\dots)$ exactly in continuous time. The time discretization does not,
so $\pi$ is reset to the equation of state at the end of every MIS step (hook
\code{frozen\_step\_end!} of \code{run\_frozen}). Within the step all substeps and stages use the
prognostic $\pi$; the drift $|\pi-\pi(\text{state})|$ before the reset is measured in the tests below.

\section{Dry air}
With $\Theta=\rho\theta$, $p=p_0(R\Theta/p_0)^{\gamma}$ and the Exner pressure
\[
\pi=\Big(\frac{p}{p_0}\Big)^{R/c_p}=\Big(\frac{R\Theta}{p_0}\Big)^{\gamma-1},\qquad
\nabla p=c_p\,\Theta\,\nabla\pi\quad\text{(exact)},\qquad
\partial_t\pi=c\,\partial_t\Theta,\quad c=\frac{\mathrm d\pi}{\mathrm d\Theta}=(\gamma-1)\frac{\pi}{\Theta}.
\]
The extended system for $(\rho,U,W,\Theta,\pi)$, $M=(U,W)$:
\begin{align*}
\partial_t\rho+\nabla\cdot M&=0, &
\partial_t\Theta+\nabla\cdot(\theta M)&=0, &
\partial_t\pi+c\,\nabla\cdot(\theta M)&=0,\\
\partial_t U+\nabla\cdot(Mu)+c_p\Theta\,\partial_x\pi&=0, &
\partial_t W+\nabla\cdot(Mw)+c_p\Theta\,\partial_z\pi+\rho g&=0.
\end{align*}

\paragraph{Frozen form.} $N(q)=F(q,q)$ with the fast part linear in the evolving $z$ and the
coefficients frozen at the stage value $y$:
\begin{align*}
\rho_t&=-\delta_x U-\delta_z W, &
\Theta_t&=-\delta_x(\theta_x U)-\delta_z(\theta_z W), &
\pi_t&=-\delta_x(\kappa_x U)-\delta_z(\kappa_z W),\\
U_t&=-c_p\Theta_x\,\delta_x\pi, &
W_t&=-c_p\Theta_z\,\delta_z\pi-g\bar\rho, &&
\end{align*}
slow: the momentum advection (shift technique). $\theta_x,\theta_z$ are the upwind face values of
$\theta$, $\Theta_x,\Theta_z$ the face means of $\Theta$, and $\kappa$ the product $c\,\theta_f$ of a cell
with its face (two coefficients per face, one for each neighbour). Products of frozen quantities are
stored as coefficients of their own, because the MIS stage combination $\sum_j w_jC(Y_j)$ must be
linear in $C$. In $N$ the base operator uses the same pressure gradient,
\code{FV\_PRESSURE\_FORM = :exner}: $c_p\Theta_f\,\delta\pi(\Theta)$.

\paragraph{HEVI.} The vertical part is linear in $z$; eliminating $\rho,\Theta,\pi$ gives per column the
tridiagonal system ($r=(c_{\rm CN}/\Delta z)^2$, $g_g=c_{\rm CN}^2g/(2\Delta z)$, $P_k=c_p\Theta_{z,k}$)
\begin{multline*}
W_{k-1}\big[-rP_k\kappa^{T}_{k-1}+g_g\big]+W_k\big[1+rP_k(\kappa^{T}_k+\kappa^{B}_k)\big]+W_{k+1}\big[-rP_k\kappa^{B}_{k+1}-g_g\big]\\
=R^W_k-\tfrac{c_{\rm CN}}{\Delta z}P_k(R^\pi_k-R^\pi_{k-1})-\tfrac{c_{\rm CN}g}{2}(R^\rho_{k-1}+R^\rho_k),
\end{multline*}
$\kappa^B,\kappa^T$ the coefficients of the cell below/above the face; the step is Crank--Nicolson.

\paragraph{Balance.} For an isentropic atmosphere $\pi$ is linear in $z$ and the analytic state satisfies
the discrete balance $c_p\Theta_f\,\delta_z\pi+g\bar\rho=0$ to round-off, without a background state
(residual $6\cdot10^{-13}$ against $8\cdot10^{-5}$ for the product form $(\mathrm dp/\mathrm d\Theta)\nabla\Theta$).
For a stratified atmosphere the initial column is balanced discretely in the pressure form of the
model ($\rho_k$ from the $W$ equation at rest).

\section{Moist air}
Formulation (I) of the poster, reversible thermodynamics as in Bryan and Fritsch (2002):
prognostic $\rho$, $M$, $\Theta=\rho\theta_\rho$, $V=\rho_v=\rho q_v$ and $C=\rho_c=\rho q_c$ (poster
eqs.~3--5; $q_t=q_v+q_c$ diagnostic),
\[
\partial_t\Theta+\nabla\cdot(\theta_\rho M)=Q_\theta,\qquad\partial_tV+\nabla\cdot(q_vM)=-S_c,\qquad
\partial_tC+\nabla\cdot(q_cM)=S_c,
\]
with the density potential temperature
$\theta_\rho$ (moist exponent) and the diagnostic pressure
\[
p=p_0\Big(\frac{R_d\,\Theta}{p_0}\Big)^{1/(1-\kappa_M)},\qquad \kappa_M=\frac{R_m}{c_{pm}},
\]
\[
R_m=R_d\rho_d+R_v\rho_v,\qquad c_{pm}=c_{pd}\rho_d+c_{pv}\rho_v+c_{pl}\rho_c,
\]
$\rho_d=\rho-\rho_v-\rho_c$; the pressure is an explicit function of $(\rho,\Theta,V,C)$, without a saturation
adjustment. Constants and $e_s(T)$ as in CGDycore.

\subsection{Phase change: Fischer--Burmeister relaxation}
Phase equilibrium is the complementarity condition $\rho_c\ge0$, $\rho_{vs}(T)-\rho_v\ge0$,
$\rho_c\,(\rho_{vs}-\rho_v)=0$, i.e.\ $\varphi_{FB}(\rho_{vs}-\rho_v,\rho_c)=0$ with the
Fischer--Burmeister function $\varphi_{FB}(a,b)=\tfrac12(a+b-\sqrt{a^2+b^2})$ (CGDycore
\code{FischerBurmeisterRhoTh}): $\varphi_{FB}<0$ supersaturated, $\varphi_{FB}>0$ subsaturated with
condensate, and then $\varphi_{FB}\le\rho_c$. Instead of an adjustment the condensation rate is an explicit
relaxation towards $\varphi_{FB}=0$,
\[
S_c=-\frac{\varphi_{FB}(\rho_{vs}(T)-\rho_v,\;\rho_c)}{\tau},\qquad \tau=\frac{\min(\Delta x,\Delta z)}{w_\tau},
\]
with the relaxation time from the grid spacing (a parcel at speed $w_\tau$ crosses a cell in $\tau$).
Evaporation never removes more than the condensate if $\Delta t\le\tau$. The latent heating is the
reversible one, at fixed $\rho$, $Q$ and $\theta_e$:
\[
Q_\theta=\Big(\frac{\partial\Theta}{\partial\rho_c}\Big)_{\rho,\,\rho_v+\rho_c,\,\theta_e}S_c,
\]
from $\mathrm d\ln\theta_e=0$ ($\mathrm dT/\mathrm dr_v$ analytic, then $\Theta(\rho,T,\rho_v,\rho_c)$ by central
differences; $\theta_e$ is conserved to second order in the condensed amount, test). Two options:
\code{micro = :slow}, $S_c$ from the stage value $y$ in the slow part $S(y)$ (needs $\tau\gtrsim\Delta t$);
\code{micro = :fast}, $S_c$ from the current fast state in every small step (explicit, needs
$\tau\gtrsim\Delta\tau$). The phase change is also a source in the pressure equation (next
subsection).

\subsection{The simplest product form of the pressure gradient}
Use the Exner function with the \emph{dry} exponent,
\[
\pi=\Big(\frac{p}{p_0}\Big)^{\kappa_d},\quad \kappa_d=\frac{R_d}{c_{pd}}
\qquad\Longrightarrow\qquad
p=p_0\,\pi^{c_{pd}/R_d},
\]
\[
\nabla p=c_{pd}\,\hat\Theta\,\nabla\pi,\qquad \hat\Theta=\frac{p}{R_d\,\pi}=\frac{p_0^{\kappa_d}\,p^{1-\kappa_d}}{R_d}.
\]
This is exact, and the coefficient $\hat\Theta$ is a function of the pressure alone: no moisture
variable enters. $\hat\Theta=\rho\,\theta_\rho^{(d)}$ is $\rho$ times the density potential temperature
with the dry exponent ($\theta_\rho^{(d)}=T_\rho/\pi$, $T_\rho=p/(\rho R_d)$), i.e.\ the momentum equation
of Bryan and Fritsch, $\nabla p/\rho=c_{pd}\theta_\rho\nabla\pi$. In terms of the prognostic $\Theta$
(moist exponent) $\hat\Theta=\Theta\,(p/p_0)^{\kappa_M-\kappa_d}$. With the moist exponent,
$\pi_M=(p/p_0)^{\kappa_M}$, the gradient contains $\ln(p/p_0)\,\nabla\kappa_M$ and is not a product.
In the frozen form $\hat\Theta_f(y)$ is the face mean of $p/(R_d\pi)$ evaluated with the prognostic
$\pi$ of $y$.

\subsection{Pressure equation}
\[
\partial_t\pi=a\,\partial_t\rho+b\,\partial_t\Theta+c\,\partial_tV+d\,\partial_tC,\qquad
(a,b,c,d)=\frac{\partial\pi}{\partial(\rho,\Theta,V,C)},
\]
so the transport part and the phase change both enter. The source of the pressure equation (the $Q_p$
of the poster) is
\[
\partial_t\pi\big|_{\rm phase}=\kappa_S\,S_c,\qquad
\kappa_S=b\Big(\frac{\partial\Theta}{\partial\rho_c}\Big)_{\theta_e}+d-c
\]
(latent heating plus the change of $R_m$ and $\kappa_M$ from vapour to liquid), frozen per cell, with
$S_c$ from $y$ (\code{:slow}) or from the fast state (\code{:fast}). For $p$ at $T=284$\,K, $p=814$\,hPa,
$\rho_v=1.1\cdot10^{-2}$, $\rho_c=9\cdot10^{-3}$\,kg/m$^3$: $\kappa_S=9.3\cdot10^5$\,Pa\,m$^3$/kg, of which
$9.1\cdot10^5$ is the heating; it agrees with the isochoric estimate $L R_m/c_{vm}=9.5\cdot10^5$ to
1.5\,\%. A condensation rate of $10^{-5}$\,kg\,m$^{-3}$\,s$^{-1}$ raises $p$ by 9\,Pa/s. Without this
term $p$ (or $\pi$) would drift from the equation of state during the step and be corrected only by the
reset.
frozen from $y$ (central differences of the explicit equation of state). With the flux divergences of
$\rho$, $\Theta$, $V$, $C$ the transport part of the $\pi$ row is $-\delta(\kappa_f F)$ with one coefficient
per face and neighbour, $\kappa_f=a+b\,\theta_f+c\,q_{v,f}+d\,q_{c,f}$ (upwind face values); the HEVI column
is the dry one with these $\kappa$, and $V$, $C$ follow from its fluxes. With $Q=0$ the moist
model reproduces the dry one (test). After every MIS step $\pi$ is recomputed from $(\rho,\Theta,V,C)$ by
the equation of state (no adjustment).

\subsection{Initial state of the moist bubble}
Saturated, neutral atmosphere $\theta_e=320$\,K, $r_t=0.02$ ($\theta_e$ of the poster, eq.~17). The
column is in \emph{discrete} hydrostatic balance with the pressure form of the model:
$p_1$ from $\mathrm dp/\mathrm dz=-g\rho$ up to the first level, then
$c_{pd}\hat\Theta_f(\pi_k-\pi_{k-1})+g\Delta z\,\bar\rho=0$ upwards (secant in $p_k$), each level
saturated with $\theta_e=320$\,K, $r_t=0.02$, split into $V$ and $C$ of phase equilibrium ($\varphi_{FB}=0$ to
round-off). The perturbation of $\theta_\rho$ by
$2\,\mathrm K\cos^2(\pi r/2r_c)$, $r_c=2$\,km, at fixed pressure and $r_t$, saturated, as in CGDycore
\code{PerturbMoistProfile}.

\section{The pressure itself as the additional variable (\code{pvar = :p})}
\paragraph{Dry.} With $p=P(\Theta)$ the pressure equation follows from the conservative $\Theta$ equation,
\[
\partial_t p=P'(\Theta)\,\partial_t\Theta=-P'(\Theta)\,\nabla\cdot(\theta M),\qquad P'(\Theta)=\frac{\gamma p}{\Theta}=\frac{c^2}{\theta},
\]
and in the continuum this \emph{is} the standard pressure equation:
\[
-P'(\Theta)\nabla\cdot(\theta M)=-c^2\nabla\cdot M-\frac{\gamma p}{\theta}\,v\cdot\nabla\theta
=-\rho c^2\,\nabla\cdot v-v\cdot\nabla p,
\]
since $\nabla p=c^2\nabla\rho+(\gamma p/\theta)\nabla\theta$. It is frozen exactly as the $\pi$ equation:
the $p$ row is $-\delta(\kappa_fF)$ with $\kappa_f=(\gamma p/\Theta)_{\rm cell}(y)\,\theta_f(y)$ (one
coefficient per face and neighbour), the momentum rows have the plain difference $-\delta p$ (no product
form needed), the fast part is linear in $z$, the HEVI column is the tridiagonal system above with
$P_k=1$, and $p$ is reset to $P(\Theta)$ after every MIS step. The frozen coefficient is the derivative
$\gamma p/\Theta$ (as in the product form); freezing the secant $p/\Theta$ instead was unstable
(\code{docs/FINDINGS.md} \S11). The full operator $N$ is the base operator with $\delta p(\Theta)$
(\code{FV\_PRESSURE\_FORM = :gradp}).

\paragraph{Moist.} $\partial_t p=a\,\partial_t\rho+b\,\partial_t\Theta+c\,\partial_tV+d\,\partial_tC$ with
$(a,b,c,d)=\partial p/\partial(\rho,\Theta,V,C)$ plus the phase source $\kappa_SS_c$:
this is the source term $Q_p$ of the poster (frozen sound speed of the non-equilibrium mixture). The code is the moist
model with $p$ in place of $\pi$ and $-\delta p$ in the momentum rows; the balanced initial column uses
$p_k-p_{k-1}+g\Delta z\,\bar\rho=0$.

\paragraph{The advective form} $\partial_tp+v\cdot\nabla p=-\rho c^2\nabla\cdot v$ can also be frozen
(fast $-(\rho c^2)(y)\nabla\cdot(M_z/\rho_y)$, slow $-v_y\cdot\nabla p_y$, $F(q,q)=N(q)$), but it is not
conservative and the pressure advection is in the slow part: $p$ then departs from $P(\Theta)$ by the
error of the slow step, and the reset after the step is a real correction, not a cosmetic one. The
conservative form above keeps the drift at round-off level times $\dt$ (tables below).

\section{Is a moist Exner pressure useful?}
A moist Exner function $\pi_M=(p/p_0)^{\kappa_M}$ is not recommended as the prognostic variable:
\begin{itemize}
\item The product form is lost: $\nabla\pi_M=\pi_M\,[\kappa_M\nabla\ln p+\ln(p/p_0)\nabla\kappa_M]$, so the
  pressure gradient needs the correction $-p\,\kappa_M^{-1}\ln(p/p_0)\,\nabla\kappa_M$, which must be
  discretized consistently with the $\nabla\pi_M$ part or the balance is lost.
\item $\pi_M$ jumps at cloud edges, $p$ does not: $\kappa_M$ depends on the water content
  ($\kappa_M=0.2759$ at the Bryan--Fritsch surface, $\kappa_d=0.2859$ for dry air). At about 5\,km the
  jump of $\pi_M$ across a cloud edge from $\Delta\kappa\approx0.01$ alone is about $5\cdot10^{-3}$, while
  the hydrostatic change of $\pi$ over one layer of 250\,m is $7.6\cdot10^{-3}$: a spurious signal of the
  size of the hydrostatic force in $\delta\pi_M$.
\item The $\pi_M$ equation gets a source whenever only $\kappa_M$ changes (phase changes,
  without motion), and the frozen coefficient would contain $\nabla\kappa_M$ of the frozen
  moisture field.
\end{itemize}
The moist $\kappa_M$ belongs in the thermodynamic variable ($\theta_\rho$, which makes
$p=p_0(R_d\rho\theta_\rho/p_0)^{1/(1-\kappa_M)}$ exact; it changes only through the latent heating $Q_\theta$), the
pressure variable keeps the dry exponent ($\pi$) or is $p$ itself; they are linked by
$\hat\Theta=\Theta(p/p_0)^{\kappa_M-\kappa_d}=p/(R_d\pi)$.

\section{Results (flat grid, MIS54 with the HEVI fast loop)}
\paragraph{Dry warm bubble} (\code{examples/frozen\_exner\_bubble.jl}): isentropic, 20\,km$\times$10\,km,
$160\times80$ cells, $\Delta\theta=2$\,K, $T=1000$\,s, errors of $\theta'$ against RK4 of the same operator.
\begin{center}\begin{tabular}{lcccc}\toprule
$\dt$ (nsub) & 2\,s (8) & 5\,s (18) & 10\,s (36) & 15\,s (54)\\\midrule
product form & $1.67\cdot10^{-2}$ & $1.90\cdot10^{-2}$ & $1.92\cdot10^{-2}$ & $2.04\cdot10^{-1}$\\
Exner $\pi$ & $1.72\cdot10^{-2}$ & $1.93\cdot10^{-2}$ & $1.98\cdot10^{-2}$ & $1.99\cdot10^{-1}$\\
pressure $p$ & $1.64\cdot10^{-2}$ & $1.89\cdot10^{-2}$ & $1.91\cdot10^{-2}$ & $2.08\cdot10^{-1}$\\\midrule
drift $\pi$ (rel.) & $5\cdot10^{-10}$ & $6\cdot10^{-9}$ & $5\cdot10^{-8}$ & $1.2\cdot10^{-7}$\\
drift $p$ (rel.) & $1\cdot10^{-9}$ & $1.2\cdot10^{-8}$ & $5\cdot10^{-8}$ & $1.1\cdot10^{-7}$\\\bottomrule
\end{tabular}\end{center}
Same solution for all three ($\theta'_{\max}=1.971$\,K, $w_{\max}=17.35$\,m/s), same accuracy and
stability (the step to $\dt=15$\,s is limited by the advective Courant number $\approx2$ of the slow
part), mass and $\Theta$ conserved to round-off. Relative drift of the additional variable before the
reset at most $10^{-7}$. Rest state of the analytic isentropic atmosphere: residual $6\cdot10^{-13}$ ($\pi$),
$6\cdot10^{-5}$ ($p$), $8\cdot10^{-5}$ (product form); only $\pi$ is linear in $z$.

\paragraph{Gravity-wave channel} (\code{examples/frozen\_exner\_channel.jl}): Skamarock--Klemp, $N=0.01$,
$U=20$\,m/s, $120\times16$ cells, discretely balanced columns. Wave error of $\theta'$ at $z=5$\,km
against RK4 (amplitude $2.3\cdot10^{-3}$\,K):
\begin{center}\begin{tabular}{lcccccccc}\toprule
$\dt$ [s] & 2 & 5 & 10 & 20 & 50 & 100 & 150 & 200\\\midrule
product form & 0.3\,\% & 2.2\,\% & 3.0\,\% & 4.6\,\% & 7.0\,\% & 4.6\,\% & 7.8\,\% & 12\,\%\\
Exner $\pi$ & 0.3\,\% & 2.2\,\% & 3.2\,\% & 4.9\,\% & 6.8\,\% & 7.9\,\% & 16\,\% & 24\,\%\\
pressure $p$ & 0.3\,\% & 2.2\,\% & 3.1\,\% & 4.6\,\% & 6.8\,\% & 4.4\,\% & 7.0\,\% & 8.5\,\%\\\bottomrule
\end{tabular}\end{center}
Equal up to $\dt=50$\,s; at large steps the pressure variable has the smallest error (8.5\,\% at
$\dt=200$\,s against 12\,\% for the product form and 24\,\% for $\pi$). Rest-state residual of the
balanced columns $\le2\cdot10^{-15}$ for all three. Long run (36000\,s, $\dt=200$\,s): MIS54, MISF53,
MIS\_FV42 unstable with all three forms, MIS\_FV53 stable with all three ($\max|w|=1.3$, $1.5$,
$1.4\cdot10^{-3}$\,m/s for product, $\pi$, $p$).

\paragraph{Long runs and the step limit of MIS54.} MIS54 in the long channel (36000\,s) with all three
forms (pressure reset after every MIS step): stable at $\dt=50$\,s ($\max|w|=1.41$, $1.34$,
$1.43\cdot10^{-3}$\,m/s for product, $\pi$, $p$), unstable at $\dt=100$ and 200\,s. The product form has
no additional variable and fails in the same way, so the limit is that of the MIS method in this channel
(vertical Courant number of the slow advection $U\dt/\Delta z=3.2$ at $\dt=100$\,s), not the pressure
variable.

\paragraph{Divergence damping by extrapolation} (\code{run\_frozen(\dots; smdiv = $\beta$)}). As in WRF,
the momentum part of the fast loop is evaluated at the extrapolated state
$\tilde z=z+\beta(z-z_{\rm prev})$, $z_{\rm prev}$ the state of the previous momentum evaluation
($z_{\rm prev}=z$ at the start of the stage); the momentum rows depend only on the thermodynamic rows,
so this extrapolates $\Theta$, $\pi$ or $p$ (HEVI: the explicit horizontal momentum steps; Verlet: the
momentum step, with gravity and sponge). It acts like a divergence damping $\sim\beta c^2\Delta\tau$.
Long channel, $\max|w|$ [m/s] after 36000\,s (``growing'': far above the stable level
$1.5\cdot10^{-3}$):
\begin{center}\begin{tabular}{llcccc}\toprule
form & method, $\dt$ & $\beta=0$ & $0.05$ & $0.1$ & $0.2$\\\midrule
product & MIS54, 100\,s & unstable & 1.2 (growing) & $5.7\cdot10^{-3}$ & $2.2\cdot10^{-3}$\\
product & MIS54, 200\,s & unstable & $1.8\cdot10^{-3}$ & $1.8\cdot10^{-3}$ & $1.9\cdot10^{-3}$\\
product & MISF53, 200\,s & unstable & $1.6\cdot10^{-3}$ & $1.6\cdot10^{-3}$ & $1.6\cdot10^{-3}$\\
$\pi$ & MIS54, 100\,s & unstable & unstable & 0.21 (growing) & $2.9\cdot10^{-3}$\\
$\pi$ & MIS54, 200\,s & unstable & $5.2\cdot10^{-3}$ & $4.9\cdot10^{-3}$ & $4.5\cdot10^{-3}$\\
$\pi$ & MISF53, 200\,s & unstable & $2.1\cdot10^{-3}$ & $2.0\cdot10^{-3}$ & $2.0\cdot10^{-3}$\\
$p$ & MIS54, 100\,s & unstable & 0.17 (growing) & $2.2\cdot10^{-3}$ & $2.0\cdot10^{-3}$\\
$p$ & MIS54, 200\,s & unstable & $1.5\cdot10^{-3}$ & $1.5\cdot10^{-3}$ & $1.6\cdot10^{-3}$\\
$p$ & MISF53, 200\,s & unstable & $1.5\cdot10^{-3}$ & $1.5\cdot10^{-3}$ & $1.5\cdot10^{-3}$\\
all & MIS\_FV42, 200\,s & unstable & unstable & unstable & unstable\\\bottomrule
\end{tabular}\end{center}
Wave error at $T=3000$\,s, $\dt=100$\,s, MIS54, $\beta=0/0.05/0.1/0.2$: product 4.6/4.8/5.4/6.8\,\%,
$\pi$ 7.9/7.6/7.3/7.0\,\%, $p$ 4.4/4.6/5.3/6.9\,\%. The extrapolation stabilizes MIS54 and MISF53 at
$\dt=200$\,s for $\beta\ge0.05$ (amplitude as with the undamped MIS\_FV53, $1.3$--$1.5\cdot10^{-3}$\,m/s) and
at $\dt=100$\,s for $\beta\ge0.1$ (product, $p$) or $0.2$ ($\pi$); it does not stabilize MIS\_FV42. The
accuracy cost of $\beta=0.1$ (the WRF value) is about one percentage point for the product form and $p$;
for $\pi$ the error even decreases. For this channel: $p$ with $\beta=0.1$ and MIS54 or MISF53.

\paragraph{Moist bubble} (\code{examples/frozen\_moist\_bubble.jl}): see Section~\ref{sec:moistres}.

\section{Moist bubble results}\label{sec:moistres}
Bryan--Fritsch moist bubble, $80\times40$ cells ($\Delta x=\Delta z=250$\,m), $T=1000$\,s, MIS54 with HEVI,
Fischer--Burmeister relaxation with $w_\tau=50$\,m/s ($\tau=5$\,s); reference RK4 of the same operator
(3574 steps). Rest-state residual of the balanced saturated column $7\cdot10^{-13}$ ($\pi$),
$7\cdot10^{-13}$ ($p$). Reference: $\theta_e'\in[-0.69,6.10]$\,K, $\max|w|=21.0$\,m/s. Maximum error of
$\theta_e'$ [K]:
\begin{center}\begin{tabular}{lccc}\toprule
$\dt$ (nsub) & 2\,s (4) & 5\,s (10) & 10\,s (19)\\\midrule
$\pi$, \code{micro = :slow} & 0.133 & 0.099 & 0.132\\
$\pi$, \code{micro = :fast} & 0.267 & 0.229 & 0.266\\
$p$, \code{micro = :slow} & 0.113 & 0.107 & 0.120\\
$p$, \code{micro = :fast} & 0.229 & 0.213 & 0.240\\\bottomrule
\end{tabular}\end{center}
$\max|w|$ 21.0 (\code{:slow}) and 21.4\,m/s (\code{:fast}); mass and total water $V+C$ conserved to round-off, total
$\Theta$ grows by $2.4\cdot10^{-4}$ (latent heating); relative drift of $\pi$ or $p$ before the reset
$\le4\cdot10^{-8}$; remaining non-equilibrium $\max|\varphi_{FB}|/\rho=1.1\cdot10^{-4}$ (the lag of the
relaxation). $\pi$ and $p$ give the same solution. \code{:slow} is stable also at $\dt=10\,\mathrm s=2\tau$
and is about twice as accurate as \code{:fast} here, and cheaper (13--14\,s against 22\,s at $\dt=10$\,s).
The errors hardly depend on $\dt$: they come from the fast loop (fixed $\Delta\tau$), as in the dry bubble.

\paragraph{Dependence on the relaxation time} ($40\times20$ cells, $\Delta x=500$\,m, $\pi$, $\dt=5$\,s,
$\max|w|$ at 1000\,s): $\tau=50$\,s ($w_\tau=10$\,m/s) 9.3\,m/s, $\tau=10$\,s 19.7, $\tau=2$\,s 23.5,
$\tau=0.5$\,s (\code{:fast}) 25.2\,m/s; \code{:slow} and \code{:fast} agree to 1\,\%. \code{:slow} fails for
$\tau=2\,\mathrm s<\dt/2$, \code{:fast} for $\tau=0.1\,\mathrm s<\Delta\tau$. With $\tau=50$\,s the lag
leaves a supersaturation of $6\cdot10^{-4}$, which also biases the diagnosed $\theta_e$ (+1.7\,K). Without
latent heating ($Q_\theta=0$, $\rho\theta_\rho$ conserved, the first version) the bubble stalls: $\Theta/\rho$
increases by 24\,K over 10\,km in the saturated neutral background, $\max|w|$ peaks at 3\,m/s after 200\,s
and decays.

\paragraph{Relaxation on the fast scale (default).} \code{micro = :fast} with
$\tau=\tau_{\rm fac}\,\Delta\tau_{\rm ac}$, $\Delta\tau_{\rm ac}=0.85\cdot0.92\min(\Delta x,\Delta z)/c_s$
($c_s=350$\,m/s, the acoustic step of the fast loop). The default is the shortest stable relaxation,
$\tau_{\rm fac}=0.25$ ($\tau=0.14$\,s at 250\,m, $\approx0.28$ of the actual small step $\Delta t/n_{\rm sub}=0.5$\,s):
$80\times40$, $\dt=2$, 5, 10\,s, $\pi$ and $p$, $\max|w|=26.8$--27.0\,m/s at 1000\,s. At $\tau_{\rm fac}=0.2$ the
run does not blow up, but the bubble is wrong ($\max|w|=8.5$--11.4\,m/s at $\dt=2$ and 10\,s, 27.1\,m/s at
5\,s); at $\tau_{\rm fac}=0.125$ it is unstable. Stable is not accurate: with $\tau<\Delta\tau$ the relaxation is
not resolved by the small step. $80\times40$, $\pi$ (the $p$ form is the same within 10\,\%), maximum
error of $\theta_e'$ against RK4 with the same $\tau$ (RK4 step $\le\tau/2$):
\begin{center}\begin{tabular}{lccccc}\toprule
$\tau_{\rm fac}$ ($\tau$) & RK4 $\max|w|$ & MIS $\max|w|$ ($\dt=5$\,s) & $\dt=2$\,s & 5\,s & 10\,s\\\midrule
0.25 (0.14\,s) & -- & 26.9 & -- & -- & --\\
0.5 (0.28\,s) & 22.2 & 25.7 & 2.2 & 1.8 & 2.2\\
1 (0.55\,s) & 22.1 & 24.2 & 1.0 & 0.86 & 0.97\\
2 (1.1\,s) & 22.0 & 23.1 & 0.52 & 0.41 & 0.50\\
4 (2.2\,s) & 21.7 & 22.3 & 0.31 & 0.27 & 0.32\\\bottomrule
\end{tabular}\end{center}
The physical solution (RK4) hardly depends on $\tau$ in this range (21.7--22.2\,m/s; 21.0 at
$\tau=5$\,s), the MIS error is the unresolved relaxation and does not decrease with $\dt$. Refining only
the small step at $\tau_{\rm fac}=0.5$ ($\dt=5$\,s, nsub $10,20,40,80$) gives $\max|w|=25.7$, 23.5, 22.5,
22.3\,m/s: convergence to RK4. Recommended: $\tau\gtrsim2\Delta\tau$ ($\tau_{\rm fac}=2$--4), i.e.\ the
shortest relaxation the fast step still resolves.

\paragraph{RK3, Lawson and ETD.} Besides MIS the frozen step takes the Wicker--Skamarock RK3
(\code{WSRK3}, MIS with $\beta=\mathrm{diag}(1/3,1/2,1)$), Lawson (\code{LRK4}, \code{LRK3}) and ETD
(\code{EB4}, \code{ETD2}) tableaus. Lawson and ETD freeze the fast part at the start of the step,
$F_f(C(q^n),\cdot)$, which is nonlinear in $Z$ (Fischer--Burmeister source with \code{micro = :fast}); the
slow forcing of stage $j$ is the remainder $\tilde S(Y_j)=N(Y_j)-F_f(C(q^n),Y_j)$. Lawson propagates the
increments with the nonlinear fast flow, ETD integrates each stage from $q^n$ with the polynomial forcing
(evaluated at the substep midpoint). The substeps are $\lceil c\,n_{\rm sub}\rceil$, so the fast step never
exceeds $\dt/n_{\rm sub}$ (with rounding, the half step of Lawson at $n_{\rm sub}=5$ had $1.25\Delta\tau$
and the relaxation at its limit broke down). Orders on the toy problem of the test suite (exact fast
stages): WSRK3 2; EB4 4, also with the fast coefficient frozen at $q^n$; ETD2 2; LRK4 and LRK3 only 2,
because $\varphi(X+\Delta S)\neq\varphi(X)+e^{\tau L}\Delta S$ for a nonlinear flow.

Bubble, $80\times40$, $\pi$, 1000\,s, maximum error of $\theta_e'$ [K] against RK4 at $\dt=2$, 5, 10\,s:
\begin{center}\small\setlength{\tabcolsep}{3pt}\begin{tabular}{lcccccc}\toprule
& MIS54 & WSRK3 & LRK4 & LRK3 & EB4 & ETD2\\\midrule
$\tau_{\rm fac}=0.25$ & 4.1/3.6/4.0 & 3.7/3.7/3.9 & 3.7/3.7/2.6 & 3.8/3.9/3.6 & 3.7/3.7/4.0 & 3.7/3.7/4.0\\
$\tau_{\rm fac}=4$ & 0.31/0.27/0.32 & 0.27/0.28/0.29 & 0.29/0.39/2.3 & 0.28/1.6/2.6 & 0.27/0.27/0.29 & 0.27/0.26/0.37\\\bottomrule
\end{tabular}\end{center}
At the default $\tau_{\rm fac}=0.25$ all schemes give the same $\max|w|\approx27$\,m/s (RK4 22.2): the error
is the unresolved relaxation in the fast step and does not depend on the slow scheme. With a resolved
relaxation ($\tau_{\rm fac}=4$, RK4 21.7\,m/s) MIS54, WSRK3, EB4 and ETD2 agree within 0.1\,K
($\max|w|$ 22.3--22.5\,m/s); Lawson degrades with $\dt$ (LRK3 $\max|w|=38.9$\,m/s at $\dt=10$\,s). The $p$
form is the same within 10\,\%. Mass and total water are conserved to round-off by all schemes. Cost per
step relative to MIS54 (fast time $\sum d_i$ resp.\ $\sum c_i$): WSRK3 0.7--0.8, ETD2 0.9, EB4 and LRK4 1.3.

\section{Open points}
\begin{itemize}
\item Relaxation time: the solution depends on $\tau$ (Section~\ref{sec:moistres}); $\tau\to0$ is the
  equilibrium limit, limited by the explicit source ($\tau\gtrsim\Delta t$ for \code{:slow},
  $\tau\gtrsim\Delta\tau$ for \code{:fast}). A linearly implicit source in the vertical column would remove
  the limit.
\item Larger errors of the dry Exner form in the channel at $\dt\ge100$\,s, while the pressure
  variable is the most accurate: in the $\pi$ form the frozen face mean $\Theta_f$ multiplies the evolving
  $\delta\pi$, in the $p$ form nothing is frozen in the momentum rows.
\item Energy formulation (II) and the combined system (III) of the poster; terrain and cut cells; 3D.
\end{itemize}
\end{document}
